\section{Алгоритм с неизменной длительностью измерительного окна, кратной периоду инжектируемого сигнала}\label{sec:Th_const}

\needspace{1\baselineskip}
Принципиальные особенности рассматриваемого алгоритма:
\begin{itemize}
    \item Длительность измерительного окна ${T_h}$ неизменна и равна наибольшему среди $K$-значений периоду инжектируемого сигнала.
    \item Длина $N$ последовательностей ${x\left[n\right]}$, ${y\left[n\right]}$, по которым вычисляются коэффициенты разложения в выражении \eqref{eqn:DFT}, неизменна для всех значений периода инжектируемого сигала ${T^{(k)}_z}$.
    \item Период дискретизации ${T_s}$ одинаков для всех значений периода инжектируемого сигнала ${T_z}$ и кратен периоду выполнения операций в системе ${T_d}$.
\end{itemize}

Математическая обработка в соответствии с \eqref{eqn:DFT} накладывает ограничение на период ${T_z}$ инжектируемого сигнала -- требуется, чтобы измерительное окно ${T_h}$ вмещало целое число ${L^{(k)} \in \mathbb{N}}$ периодов ${T_z}$.

Таким образом, в серии из $K$ изменений для каждого значения периода инжектируемого сигнала ${T^{(k)}_z}$, ${k \in \{1, \dots, K\}}$ (рисунок~\ref{fig:Th_const}), выполняется:
\begin{equation*}
    T^{(k)}_h = L^{(k)} \mathord{\cdot} T^{(k)}_z = N^{(k)} \mathord{\cdot} \underbrace{M^{(k)} \mathord{\cdot} T_d}_{T^{(k)}_s},
\end{equation*}
причем ${T^{(k)}_h = const}$.

\begin{figure}[H]
    \centering
    \includegraphics [width=1.00\linewidth] {frequencyResponse/Th_const.pdf}

    \caption{Параметры в алгоритме с неизменной длительностью измерительного окна, кратной периоду инжектируемого сигнала}
    \label{fig:Th_const}
\end{figure}

Частота инжектируемого сигнала потенциально принимает значения
\begin{equation*}
    f^{(k)}_z = L^{(k)} \cfrac{1}{T_h} = L^{(k)} \Delta f_z,
\end{equation*}
где ${\Delta f_z}$ -- шаг приращения частоты инжектируемого сигнала. Составляющей с частотой ${f^{(k)}_z}$ соответствуют коэффициенты ${X^{(k)}\left[l\right]}$, ${Y^{(k)}\left[l\right]}$ выражения \eqref{eqn:complexFreqResp} при ${l = L^{(k)}}$.

В связи с ${\Delta f_z = const}$ потенциальные значения частот инжектируемого сигнала ${f_z}$ равномерно распределены на интервале возможных значений, что \textbf{является главной отличительной особенностью алгоритма} в сравнении с алгоритмами \ref{sec:N_var}, \ref{sec:N_const}.

\needspace{1\baselineskip}
Изложенному алгоритму измерения характерно следующее:
\begin{enumerate}
    \item Максимально возможный период инжектируемого сигнала ${T_z}$ ограничен только параметром ${N_\text{max}}$.

    \item Длина $N$ каждого массива-отклика ${x\left[n\right]}$, ${y\left[n\right]}$ неизменна (не зависит от периода инжектируемого сигнала ${T_z}$).
    
    \item В случае ${N_\text{min}=3}$ максимально возможная частота ${f_z}$ ограничена сверху значением ${f_d / 3}$.

    \item Шаг приращения частоты инжектируемого сигнала ${f_z}$ (разрешающая способность): ${\Delta f_z = 1 / T_h}$.
    
    \item Период $T_z$ инжектируемого сигнала потенциально принимает значения от ${T^{min}_z = 3 T_s}$ до ${T^{max}_z = T_h}$, но при условии ${T_h / T_s \in \mathbb{N}}$.
\end{enumerate}